% !PRINTED 762
\weberpage{762}
This lemma allows us first of all to make, in the proof of Theorem (1), the simplifying assumption that the coefficients $C_1, C_2, \ldots, C_m$ are integral rational numbers.

For suppose that an equation

\weq{5}{X_1e^{x_1}+X_2e^{x_2}+\cdots = 0}

holds, with algebraic coefficients $X_1, X_2, \ldots$ and with exponents $x_1, x_2, \ldots$ different from one another. Then we may always assume that the $X_1, X_2, \ldots$ belong to one and the same algebraic field $\Omega$.

We denote variables by $u_1, u_2, \ldots$ and form the norm of the linear form $X_1u_1+X_2u_2+\cdots$

\weq{6}{N(X_1u_1+X_2u_2+\cdots) = \Phi(u_1,u_2,\ldots),}

which is an integral homogeneous function of the variables $u$ with \gesperrt{rational coefficients}, and whose degree is equal to the degree of the field $\Omega$. If in (6) we set

\[
u_1 = e^{x_1},\ u_2 = e^{x_2},\ \ldots,
\]

then every product of variables $u$ passes over into a quantity of the form $e^z$, in which $z$ is a sum of numbers $x$, and $\Phi(u_1,u_2,\ldots)$ becomes an expression of the form $C_1e^{z_1}+C_2e^{z_2}+\cdots$ whose coefficients are rational numbers. By 2., these coefficients cannot all be zero, provided the $X_1, X_2, \ldots$ do not all vanish, and they cannot all be zero even when those terms of the function $\Phi(e^{x_1},e^{x_2},\ldots)$ which are equal to one another are collected into a single term $Ce^z$. But if now equation (5) holds, then $\Phi(e^{x_1},e^{x_2},\ldots)=0$ as well, and consequently

\[
C_1e^{z_1}+C_2e^{z_2}+\cdots = 0.
\]

If there are fractional numbers among the $C_1, C_2, \ldots$, we can convert them into integers by multiplying the whole equation by the common denominator.

3. \gesperrt{It therefore only remains to be proved that equation (1) can hold for no system of integral rational numbers $C_1, C_2, \ldots, C_m$ which do not all vanish.}

Accordingly we now assume that an equation

\weq{7}{A_1e^{x_1}+A_2e^{x_2}+\cdots = 0,}

holds whose coefficients $A_1, A_2, \ldots$ are integral rational numbers,
%CONT>>
